%%%%%%%%%%%%%%%%%%%%%%%%%%%%%%%%%%%%%%%%%%%%%%%%%%%%%%%%%%%%%%%%%%%%%%%%%%%%%%%%%%%%%%%%%%%%%%%%%%%%%%%%%%%%%%%%%%%%%%%%%%%%%%%%%%%%%%%%%%%%%%%%%%%%%%%%%%%
% This is just an example/guide for you to refer to when submitting manuscripts to Frontiers, it is not mandatory to use Frontiers .cls files nor frontiers.tex  %
% This will only generate the Manuscript, the final article will be typeset by Frontiers after acceptance.   
%                                              %
%                                                                                                                                                         %
% When submitting your files, remember to upload this *tex file, the pdf generated with it, the *bib file (if bibliography is not within the *tex) and all the figures.
%%%%%%%%%%%%%%%%%%%%%%%%%%%%%%%%%%%%%%%%%%%%%%%%%%%%%%%%%%%%%%%%%%%%%%%%%%%%%%%%%%%%%%%%%%%%%%%%%%%%%%%%%%%%%%%%%%%%%%%%%%%%%%%%%%%%%%%%%%%%%%%%%%%%%%%%%%%

%%% Version 3.4 Generated 2018/06/15 %%%
%%% You will need to have the following packages installed: datetime, fmtcount, etoolbox, fcprefix, which are normally inlcuded in WinEdt. %%%
%%% In http://www.ctan.org/ you can find the packages and how to install them, if necessary. %%%
%%%  NB logo1.jpg is required in the path in order to correctly compile front page header %%%

\documentclass[utf8]{frontiersSCNS} % for Science, Engineering and Humanities and Social Sciences articles

\usepackage{url,hyperref,lineno,listings,microtype,subcaption}

% Automatic formatting of SI units
\usepackage[binary-units]{siunitx}

\usepackage[onehalfspacing]{setspace}

% Required for 'straight' quotes in code listings
\usepackage[T1]{fontenc}

% Visible TODO notes
\newcommand{\todo}[1]{\textbf{\textsc{\textcolor{red}{(TODO: #1)}}}}

\linenumbers


% Leave a blank line between paragraphs instead of using \\


\def\keyFont{\fontsize{8}{11}\helveticabold }
\def\firstAuthorLast{Sample {et~al.}} %use et al only if is more than 1 author
\def\Authors{James C Knight\,$^{1,*}$, Thomas Nowotny\,$^{1}$}
% Affiliations should be keyed to the author's name with superscript numbers and be listed as follows: Laboratory, Institute, Department, Organization, City, State abbreviation (USA, Canada, Australia), and Country (without detailed address information such as city zip codes or street names).
% If one of the authors has a change of address, list the new address below the correspondence details using a superscript symbol and use the same symbol to indicate the author in the author list.
\def\Address{$^{1}$Centre for Computational Neuroscience and Robotics, School of Engineering and Informatics, University of Sussex, Brighton, United Kingdom }
% The Corresponding Author should be marked with an asterisk
% Provide the exact contact address (this time including street name and city zip code) and email of the corresponding author
\def\corrAuthor{James C Knight}

\def\corrEmail{J.C.Knight@sussex.ac.uk}


\begin{document}
\onecolumn
\firstpage{1}

\title[Running Title]{Article Title} 

\author[\firstAuthorLast ]{\Authors} %This field will be automatically populated
\address{} %This field will be automatically populated
\correspondance{} %This field will be automatically populated

\extraAuth{}% If there are more than 1 corresponding author, comment this line and uncomment the next one.
%\extraAuth{corresponding Author2 \\ Laboratory X2, Institute X2, Department X2, Organization X2, Street X2, City X2 , State XX2 (only USA, Canada and Australia), Zip Code2, X2 Country X2, email2@uni2.edu}


\maketitle


\begin{abstract}

%%% Leave the Abstract empty if your article does not require one, please see the Summary Table for full details.
\section{}
Inspired by the distributed nature of memory and computation in the brain, neuromorphic systems aim to re

However, while these systems may be the ultimate platform for \textit{deploying} spiking neural networks, their very distributed nature makes them unwieldy tools for \textit{developing} SNN models.
Instead development of SNN models tends to be performed on computers with standard Von Neumann architectures and, once models scale above a certain size, clusters of Von Neumann architecture machines.

In this paper we re-implement two large-scale point neuron network models with levels of connectivity and sparseness matching those in the neocortex using our GeNN code generator.
We then verify the correctness of our GPU simulations before comparing the performance and energy usage against published data for traditional CPU-based HPC as well as neuromorphic hardware.

\tiny
 \keyFont{ \section{Keywords:} keyword, keyword, keyword, keyword, keyword, keyword, keyword, keyword} %All article types: you may provide up to 8 keywords; at least 5 are mandatory.
\end{abstract}

\section{Introduction}

\begin{itemize}
    \item Neuromorphic hardware families - limitations, scales etc
    \item HPC - Neuron, NEST etc, recent benchmarks
    \item GPU hardware introduction
    \item Other GPU simulators
\end{itemize}

\section{Material and Methods}
\subsection{GeNN}
\begin{itemize}
    \item GeNN code generator - while CUDA is relatively simple, efficiently parallelising SNN code is non-trivial
    \item Does not rely on any sort of domain-specific language - models are described by snippets of serial C-like code
    \item Walk through example of a neuron model
\end{itemize}

GeNN neuron models are defined by writing a class defining the model parameters and snippets of C-like code that describe how it should be simulated.
For example the following `LIF' class describes a very simple leaky integrate-and-fire neuron with normalised units, solved algebraically:
%
\lstinputlisting[language=C++,showstringspaces=false,basicstyle=\tiny,numbers=left,upquote=true,firstline=2,lastline=15]{code_snippets.cc}
%
The \lstinline{DECLARE_MODEL} and \lstinline{IMPLEMENT_MODEL} macros define boilerplate code used for subsequently defining parameters and initial model states.
\todo{sim code, parameters, values, derived parameters}
Macros are used to provide some `syntactic sugar' and type safety but are entirely optional.
GeNN uses this same approach for defining synapse models with the option to define code snippets for continuous time-driven updates as well as event-driven updates in response to pre and postsynaptic spikes.
Once r
%
\lstinputlisting[language=C++,showstringspaces=false,basicstyle=\tiny,numbers=left,upquote=true,firstline=17,lastline=20]{code_snippets.cc}
%
This also demonstrates how the same approach used for defining models can be used to initialise variables.
Here the membrane voltage for our \num{1000} LIF neurons is sampled from the uniform distribution using a `variable initialisation snippet' built into GeNN which allows the same parallelisation strategies used for simulating models on the GPU to be applied to initialising them.
This approach is advantageous both as it removes the need to transfer the model state from the CPU to the GPU
While the PCI express bus typically used to connect the GPU to the host system is relatively fast (up to \SI{16}{\giga\byte\per\second}), initialising millions of synaptic weights on the CPU and then uploading them at the start of the simulation can still considerably slow down the time it takes to iterate the model.

Models can be implemented using the C++ interface described in the preceding section and the resultant simulation code can be linked against 
\todo{Brian2GeNN, SpineML, PyNN}
\subsection{The cell-type specific cortical microcircuit}
sfa
\subsection{Balanced random network with power-law STDP}
The type of one-shot, online learning observed in nature is one of the many properties that make spiking neural networks attractive when compared to the approaches typically employed in training artificial neural networks.
\todo{link between learning and synaptic plasticity}
However, adding synaptic plasticity to spiking neural network simulations typically increases the computational cost of simulating them significantly. 
\citet{Morrison2007} reported that adding plasticity to their simulations slowed them down by ``a factor of less than 10'' and \citet{Knight2016b} found that, in the best case, simple STDP plasticity reduced the performance of the SpiNNaker neuromorphic system by approximately $6\times$.

\section{Results}
\subsection{Model }
\subsection{Correctness}
jeldh

\subsection{Performance}
sss

\subsection{Power and energy}
sda

\section{Discussion}
\subsection{Neurorobotics}
afa

\subsection{Interactive simulation}
asfs

\subsection{Future GPU performance}
Since the end of ``the free lunch'' -- the point at which it became impractical to continue building ever faster and more complex single-core CPUs -- the only way processor designers have been able to extract more performance from the increasing number of transistors offered by newer fabrication technologies was by build multi-core CPUs. \todo{citations}
However parallel CPU programming paradigms such as multithreading are not well-suited to the type of \textit{fine-grained} parallelism required to harness hundreds of cores.\todo{citations}

The SIMT programming model architecture used by NVIDIA GPUs is designed explicitly for fine-grained parallelism and high throughput.
However spiking neural networks are not especially well-suited

\begin{itemize}
    \item HBM3 - further doubling of memory bandwidth
    \item FP16 - half memory bandwidth, double arithmetic throughput
\end{itemize}

\section*{Conflict of Interest Statement}

The authors declare that the research was conducted in the absence of any commercial or financial relationships that could be construed as a potential conflict of interest.

\section*{Author Contributions}

The Author Contributions section is mandatory for all articles, including articles by sole authors. If an appropriate statement is not provided on submission, a standard one will be inserted during the production process. The Author Contributions statement must describe the contributions of individual authors referred to by their initials and, in doing so, all authors agree to be accountable for the content of the work. Please see  \href{http://home.frontiersin.org/about/author-guidelines#AuthorandContributors}{here} for full authorship criteria.

\section*{Funding}
Details of all funding sources should be provided, including grant numbers if applicable. Please ensure to add all necessary funding information, as after publication this is no longer possible.

\section*{Acknowledgments}
This is a short text to acknowledge the contributions of specific colleagues, institutions, or agencies that aided the efforts of the authors.

\section*{Supplemental Data}
 \href{http://home.frontiersin.org/about/author-guidelines#SupplementaryMaterial}{Supplementary Material} should be uploaded separately on submission, if there are Supplementary Figures, please include the caption in the same file as the figure. LaTeX Supplementary Material templates can be found in the Frontiers LaTeX folder.

\section*{Data Availability Statement}
The datasets [GENERATED/ANALYZED] for this study can be found in the [NAME OF REPOSITORY] [LINK].
% Please see the availability of data guidelines for more information, at https://www.frontiersin.org/about/author-guidelines#AvailabilityofData

\bibliographystyle{frontiersinSCNS_ENG_HUMS} % for Science, Engineering and Humanities and Social Sciences articles, for Humanities and Social Sciences articles please include page numbers in the in-text citations
%\bibliographystyle{frontiersinHLTH&FPHY} % for Health, Physics and Mathematics articles
\bibliography{test}

%%% Make sure to upload the bib file along with the tex file and PDF
%%% Please see the test.bib file for some examples of references

\section*{Figure captions}

%%% Please be aware that for original research articles we only permit a combined number of 15 figures and tables, one figure with multiple subfigures will count as only one figure.
%%% Use this if adding the figures directly in the mansucript, if so, please remember to also upload the files when submitting your article
%%% There is no need for adding the file termination, as long as you indicate where the file is saved. In the examples below the files (logo1.eps and logos.eps) are in the Frontiers LaTeX folder
%%% If using *.tif files convert them to .jpg or .png
%%%  NB logo1.eps is required in the path in order to correctly compile front page header %%%

\begin{figure}[h!]
\begin{center}
\includegraphics[width=10cm]{logo1}% This is a *.eps file
\end{center}
\caption{ Enter the caption for your figure here.  Repeat as  necessary for each of your figures}\label{fig:1}
\end{figure}


\begin{figure}[h!]
\begin{center}
\includegraphics[width=15cm]{logos}
\end{center}
\caption{This is a figure with sub figures, \textbf{(A)} is one logo, \textbf{(B)} is a different logo.}\label{fig:2}
\end{figure}

%%% If you are submitting a figure with subfigures please combine these into one image file with part labels integrated.
%%% If you don't add the figures in the LaTeX files, please upload them when submitting the article.
%%% Frontiers will add the figures at the end of the provisional pdf automatically
%%% The use of LaTeX coding to draw Diagrams/Figures/Structures should be avoided. They should be external callouts including graphics.

\end{document}
